%! Rational Functions and Zeros — Unit 1, Lesson 1.8 — Student Packet Cover Sheet
\documentclass[10pt]{article}
\usepackage{precalculus-article}
\usepackage{precalculus-boxes}
\usepackage{ltablex}
\keepXColumns

\begin{document}

% Full-bleed navy banner
\begin{tikzpicture}[remember picture, overlay]
  \fill[navy] (current page.north west)
    rectangle ([yshift=-0.9in]current page.north east);
\end{tikzpicture}
\vspace{-0.8in}

\begin{center}
  {\color{white}\textbf{\LARGE Precalculus}}\\[4pt]
  {\color{skymid}\large\textbf{Unit 1: Polynomial and Rational Functions}}\\[6pt]
  {\color{white}\textbf{\large Lesson 1.8 \quad Rational Functions and Zeros}}
\end{center}

\vspace{0.22in}
\namedateperiod
\vspace{0.18in}

\begin{learningtargetbox}
\small
\begin{enumerate}[leftmargin=1.5em, itemsep=5pt]
  \item I can find the real zeros of a rational function by setting the numerator equal to zero
        and verifying that those values are in the domain (i.e., do not make the denominator zero).
  \item I can solve a rational inequality by building a sign chart with critical values from both
        the numerator and the denominator, and writing the solution in interval notation.
\end{enumerate}
\end{learningtargetbox}

\vspace{0.14in}

\begin{tocbox}
\small
\renewcommand{\arraystretch}{1.5}
\begin{tabularx}{\linewidth}{@{} c l X r @{}}
\rowcolor{sky}
\textbf{\#} & \textbf{Component} & \textbf{Description} & \textbf{Score} \\
\midrule
1 & Warm-Up        & Spiral review: horizontal asymptotes and factoring      & \blank{1.2cm} \\
2 & Group Activity & Break-even analysis using zeros and sign charts          & \blank{1.2cm} \\
3 & Exit Ticket    & Independent check on zeros and a domain-restriction error & \blank{1.2cm} \\
4 & Homework       & Practice and AP-style extension                          & \blank{1.2cm} \\
\midrule
  & \multicolumn{2}{r}{\textbf{Total}} & \blank{1.2cm} \\
\end{tabularx}
\end{tocbox}

\end{document}
